\begin{figure}[!ht]
  \centering
  \begin{tikzpicture}[alt={The same native Memory interval after reconstruction. Six outlined bytes are 0f 19 23 2d 37 41, representing 15, 25, 35, 45, 55 and 65. Selectors zero and two retain their positions.}]
    \node[anchor=north west,inner sep=0,font=\small\bfseries] (title) at (0,0) {The same bytes now represent RGB samples};
    \node[anchor=north west,inner sep=0] (a) at ([yshift=-2mm]title.south west) {\MARSRevisionCrop{revision/captures/png/png-reconstructed.png}{0.99\linewidth}{0.03125}{0.33888889}{0.58203125}{0.60888889}};
    \begin{scope}[shift={(a.south west)},x={($(a.south east)-(a.south west)$)},y={($(a.north west)-(a.south west)$)}]
      \draw[black,line width=.7pt] (.5825,.01) rectangle (.869375,.467);
      \draw[black,dashed,line width=.7pt] (.20,.01) rectangle (.2478125,.467);
      \draw[black,dashed,line width=.7pt] (.5346875,.01) rectangle (.5825,.467);
    \end{scope}
  \end{tikzpicture}
  \caption{After reconstruction, the outlined samples are decimal 15, 25, 35, 45, 55 and 65. The selectors remain zero and two. The complete buffer retains its thirty-five-byte extent, although its sample bytes have changed meaning.}
  \label{fig:revision-png-native-reconstructed}
\end{figure}
